% C2O — Close-to-Open: a cross-sectional, dollar-neutral overnight strategy on US large caps.
% Machine Learning in Finance — group coursework report.
% Compile with pdflatex (twice for references). No external bib needed.
\documentclass[11pt,a4paper]{article}

\usepackage[margin=2.4cm]{geometry}
\usepackage{amsmath,amssymb}
\usepackage{booktabs}
\usepackage{graphicx}
\graphicspath{{figures/}}
\usepackage{xcolor}
\usepackage{enumitem}
\usepackage{hyperref}
\hypersetup{colorlinks=true,linkcolor=blue!50!black,urlcolor=blue!50!black,citecolor=blue!50!black}
\usepackage{titlesec}
\titleformat{\section}{\large\bfseries}{\thesection}{0.6em}{}
\titleformat{\subsection}{\normalsize\bfseries}{\thesubsection}{0.6em}{}
\setlist{nosep,leftmargin=1.4em}
\newcommand{\code}[1]{\texttt{#1}}

% ===========================================================================================
% Headline numbers — SINGLE SOURCE OF TRUTH (reproduced by the code under data/outputs/<run_id>/).
% Update these from the run's tables/*.csv; the prose references the macros only.
% ===========================================================================================
\newcommand{\reconResid}{4.4\times10^{-16}}
% --- alpha / sleeves (OOS 2015--2024 daily Spearman IC) ---
\newcommand{\icRidge}{0.016}
\newcommand{\icHGB}{0.024}
\newcommand{\icEns}{0.022}
\newcommand{\icEnsT}{7.2}
\newcommand{\icFlow}{0.014}
\newcommand{\icFlowT}{6.8}
\newcommand{\icComb}{0.024}
\newcommand{\sleeveCorr}{0.11}
\newcommand{\sleeveRetCorr}{0.20}
% --- DL benchmark (OOS IC) ---
\newcommand{\dlHGB}{0.024}
\newcommand{\dlMLP}{0.020}
\newcommand{\dlMLPseq}{0.020}
% --- v1 baseline and v2 headline / frontier net Sharpe (net of full Section 6.3 costs) ---
\newcommand{\netVone}{0.38}                  % v1 headline: ensemble, 2\% tails, @250M
\newcommand{\headSmall}{0.32}                % HEADLINE ensemble 1.25\% @50M
\newcommand{\headBench}{0.71}                % HEADLINE @250M (peak)
\newcommand{\headLarge}{0.58}                % HEADLINE @1B
\newcommand{\headPosYears}{8/10}
\newcommand{\grossHead}{1.81}                % gross Sharpe of headline @250M
\newcommand{\frontSmall}{0.74}              % FRONTIER HGB 1\% @50M
\newcommand{\frontBench}{0.80}             % FRONTIER @250M (peak; best single number)
\newcommand{\frontLarge}{0.05}             % FRONTIER @1B (collapses: 1\% book thins to ~5 names)
% --- structural ---
\newcommand{\edgeCorr}{-0.003}            % corr(predicted day edge, realised net return): no day-timing
\newcommand{\slipLoss}{0.80}              % Sharpe lost to slippage at the headline

\title{\textbf{C2O --- Close-to-Open}\\[2pt]
\large A cross-sectional, dollar-neutral overnight long/short strategy on US large-cap equities}
\author{Daniel Leung}
\date{Development window: 2010--2024 \quad(held-out evaluation: 2025--2026, single \code{cutoff} parameter)}

\begin{document}
\maketitle

\begin{abstract}
\noindent We build a daily, dollar-neutral, US large-cap long/short strategy that holds each position for
exactly one overnight session (closing-auction entry, next opening-auction exit), point-in-time throughout
and keyed off a single anti-leakage \code{cutoff}. The cross-sectional overnight alpha is genuinely present
out-of-sample (ensemble mean daily IC $\approx\icEns$, $t\approx\icEnsT$). The decisive constraint is
\emph{economic}: the mandated $4$\,bps round trip is charged on a \emph{forced} nightly liquidation of the
\emph{gross} book, and dollar-neutrality strips the universal overnight drift, so net profitability needs the
cross-sectional half-spread to beat a fixed per-gross cost. We treat the search for net Sharpe as a controlled
experiment. We (i) add an orthogonal fundamental-flow sleeve (analyst revisions / earnings surprise; IC
$\approx\icFlow$, rank-correlation only $\sleeveCorr$ to the price sleeve), (ii) benchmark a neural model
against the gradient-boosted tree, and (iii) test cost-aware per-name selection, edge-proportional gross
scaling, sector/beta neutralization, vol-targeting and inverse-vol weighting. The disciplined finding,
documented with a full ablation, is that \textbf{tail concentration is the only lever that lifts net Sharpe};
the sophisticated overlays do not beat it, for reasons we make precise (the cost is on gross, so per-name
selection cannot reduce it; and predicted day-quality has $\approx\edgeCorr$ correlation with realised
returns, so the good nights cannot be timed). Correctly locating the broad concentration optimum
($\approx1.25\%$ tails) roughly \emph{doubles} the v1 net Sharpe from $\netVone$ to a peak of $+\headBench$
(headline) and $+\frontBench$ (higher-risk frontier) at \$250M, positive in $\headPosYears$ years. We did not
reach a Sharpe of one and do not claim to: we show, with the cost arithmetic, that net Sharpe $\approx0.7$--$0.8$
is the honest ceiling of this mandate, and that crossing one would require relaxing a constraint of the mandate
itself. We report these numbers rather than inflated ones, and document every lever we tried --- including those
that failed.
\end{abstract}

\section*{Reproducibility (read first)}
Every number and figure is produced by one command, \code{python -m c2o.main}, which writes a self-describing
\code{data/outputs/<run\_id>/} (tables, figures, QuantStats tear-sheet, \code{manifest.json} with config + git
SHA). The neural benchmark is a separate, off-path experiment (\code{python -m c2o.dl\_alpha}). Construction
variants were explored in a cached lab (\code{tools/experiment*.py}) and the winners promoted to
\code{config/default.yaml}; nothing in the headline is hand-set in code. Fast verification: \code{pytest} (26
unit tests $<5$\,s). The development window is strictly 2010--2024; moving \code{window.cutoff} evaluates
2025--2026 with no code change. Out-of-sample results begin 2015 (2010--2014 is the walk-forward burn-in).

% =====================================================================================
\section{The overnight anomaly --- economic context and thesis}
Before any modelling it is worth being explicit about \emph{why} an overnight strategy could work and what kind
of edge we are hunting, because that prior shaped every later design choice.

\paragraph{Why overnight returns are special.} The overnight session is not a quiet continuation of the day. It
is when most value-relevant information arrives --- earnings (the large majority of US large-caps report before
the open or after the close), analyst actions, macro prints and overseas trading --- yet it is a window in which
participants cannot trade continuously. Three mechanisms follow. (i) An \emph{overnight risk premium}: holders
of inventory across the close demand compensation for gap risk they cannot hedge intraday, which appears as a
systematic positive close-to-open drift (the brief's Section~1 stylised fact, which we reproduce). (ii)
\emph{Information and inventory effects}: the opening auction clears orders accumulated overnight, so the
open-vs-prior-close move is unusually rich in news and dealer-rebalancing flow. (iii) A \emph{cross-sectional}
component: names that gapped, that had unusual late volume, or that sit at liquidity extremes revert or continue
overnight in partly predictable ways. The first mechanism is a market-wide \emph{level} effect; the third is a
\emph{cross-sectional} effect. This distinction is the fulcrum of the whole report.

\paragraph{Our thesis, stated up front.} A dollar-neutral book cannot earn the level effect --- it is common to
longs and shorts and cancels exactly. It can only earn the cross-sectional component. So the enterprise rests on
whether the cross-sectional overnight signal is (a) real out-of-sample and (b) large enough to clear the mandated
frictions. We will show (a) is comfortably true (the signal is strong and persistent) but (b) is the binding
constraint --- and that this one fact, a strong gross edge throttled by a structural cost, explains every result
that follows, including which sophisticated ideas help and which do not.

\paragraph{What we deliberately do \emph{not} assume.} We do not assume the level drift can be harvested (it
cannot, under dollar-neutrality); we do not assume costs can be engineered away (they are mandated and, as we
show, structural); and we do not assume a bigger model is the answer (it is not --- the gross signal is already
strong). Naming these non-assumptions early is what kept the project out of the three commonest dead ends, which
we revisit in \S\ref{sec:overfit}.

% =====================================================================================
\section{Step 1 --- Building the daily panel (brief §2)}

\subsection{Return objects and reconciliation}
With $\mathrm{Open}_t,\mathrm{Close}_t$ the split/dividend-adjusted auction prices, we define
$r_{\mathrm{ON},t}=\mathrm{Open}_t/\mathrm{Close}_{t-1}-1$,
$r_{\mathrm{ID},t}=\mathrm{Close}_t/\mathrm{Open}_t-1$,
$r_{\mathrm{CC},t}=\mathrm{Close}_t/\mathrm{Close}_{t-1}-1$. We place raw OHLC on the adjusted-close scale via
a single factor $\mathrm{adjusted\_close}/\mathrm{close}$, so open and close are corporate-action-consistent
\emph{by construction}; hence $(1+r_{\mathrm{ON}})(1+r_{\mathrm{ID}})-1\equiv r_{\mathrm{CC}}$. Empirically the
identity holds to machine precision: the maximum absolute residual over all valid stock-days is
$\reconResid$, far inside our tolerance $10^{-8}$; zero stock-days fail. Adjusting both legs by the same factor
removes the class of bug (inconsistent split handling between open and close) the brief warns about; the
residual then only tests floating-point arithmetic, which we state plainly rather than over-claim.

\paragraph{Q2.1 (data source, window, survivorship).} Source: the supplied daily public-data files
(\code{prices}, \code{earnings\_calendar}, \code{earnings\_transfo}, \code{short\_interest\_transfo},
\code{cheapness\_scores}, \code{gics\_info}, \code{regime}, \code{sp500\_tr}). Window: 2010--2024 for modelling;
pre-2010 prices retained only for lookback features. Survivorship: the panel keeps every instrument through its
last observed date and selects the universe point-in-time at each year-start (below), so delisted names remain
in the historical panel --- no retroactive dropping.

\paragraph{Q2.2 (corporate-action reconciliation, tolerance, failures).} Adjusted-scale returns, tolerance
$10^{-8}$, max residual $\reconResid$, $0\%$ of stock-days fail. We additionally flag and exclude from
\emph{ranking} severe raw-price scale anomalies (\code{close}$>$10{,}000 with \code{adjusted\_close}$<$1{,}000)
so a single corrupt print cannot distort the market-cap sort.

\paragraph{Q2.3 (BMO/AMC encoding, verified by hand).} An after-market-close (AMC) print on day $D$ is unknown
to a 15:50 trader on $D$; a before-market-open (BMO) print on $D$ is known. For the \emph{event-risk exclusion}
we flag the decision date whose overnight session crosses the release: AMC on $D\Rightarrow$ exclude $D$; BMO on
$D\Rightarrow$ exclude $D{-}1$; missing timing excludes both, conservatively. \emph{Hand check} (unit-tested in
\code{tests/test\_panel.py}): a synthetic AMC event on 2020-01-06 flags 2020-01-06; a BMO event on 2020-01-06
flags the prior session 2020-01-03.

\paragraph{Q2.4 (short-interest construction and lag).} The supplied series already embeds FINRA's
$\sim$8-day publication lag plus the required 2-day vendor lag. On top we build a daily column by an
\emph{as-of-previous-trading-day} merge: the value usable for a 15:50 decision on day $t$ is the most recent
snapshot available on or before $t{-}1$, forward-filled between releases. No same-day short-interest can leak.

\paragraph{Q2.5 (universe evolution).} On the first trading day of year $Y$ we rank all instruments by market
cap at the last trading day of $Y{-}1$ (requiring $\ge$12 months of history), take the top $1{,}000$ and
\emph{freeze} that list for the year; delistings leave on their last date. The selection yields exactly
$1{,}000$ names each year; mid-year exits are recorded (\code{step1\_universe\_counts}).

\paragraph{Q2.6 (stylised fact).} The equal-weighted overnight stream of the frozen universe dominates the
close-to-close stream over 2010--2024, while the intraday stream is materially weaker --- reproducing the
brief's Section~1 picture. This positive overnight \emph{drift} is precisely what dollar-neutrality later
cancels, central to our diagnosis (\S6).

% =====================================================================================
\section{Step 2 --- A capacity-aware universe (brief §3)}

\paragraph{Q3.1 (thresholds, justified).} Per-share price floor \$5; year-start market-cap floor \$2B;
annualised realised-volatility band $[5\%,120\%]$; 20-day windows for ADV, volatility and the Roll spread; an
earnings exclusion window of the event-risk overnight $+1$ trading day. ADV20 and volatility are shifted one
day and computed on the full price history, so early-2010 rows use valid 2009 lookback without look-ahead.

\paragraph{Q3.2 (participation cap and implied slippage).} We cap any single position at $5\%$ of trailing
ADV20. Under the square-root convention $\text{impact}\approx k\,\sigma_{\text{daily}}\sqrt{f}\times10^4$ with
$k=0.7$, $f=0.05$, the implied one-way slippage at the cap is single-digit bps for a typical large cap ---
consistent with the fixed $1.5$\,bps/leg auction slippage used for headline costs (\S6).

\paragraph{Q3.3--Q3.4 (eligible-set evolution by AUM; binding-constraint distribution).} The eligibility table
labels each (stock, date, AUM) by the single binding reason (\code{OK}, \code{ADV\_FAIL}, \code{MCAP\_FAIL},
\code{PRICE\_FAIL}, \code{VOL\_FAIL}, \code{EARN\_WINDOW}). Median daily \code{OK} names fall from $\approx852$
at \$50M to $\approx714$ at \$250M to $\approx379$ at \$1B --- the expected capacity contraction, driven by
\code{ADV\_FAIL}. The volatility floor bites visibly in 2020\,Q1. \emph{Design note}: Step~2 uses the brief's
$200$-name \emph{planning} size for the ADV test; the actual book is sized exactly by the Step~5 water-fill, so
the residual capacity shows up honestly as gross-utilisation $<1$ (\S6).

% =====================================================================================
\section{Step 3 --- Filtering the short leg for borrow (brief §4)}

\paragraph{Q4.1 (hard-to-borrow proxy).} From lagged short-interest features --- short-interest ratio
\code{dsi}, days-to-cover \code{dtcn}, change \code{ddtcn} --- we combine absolute thresholds
($\code{dsi}\ge10\%/20\%$, $\code{dtcn}\ge7.5/12.5$, $\code{ddtcn}\ge2/5$) with daily cross-sectional
$90^{\text{th}}/98^{\text{th}}$ percentile flags and a size/liquidity context flag. An additive stress score
maps to tiers A (general collateral), B (score$\ge2$), C (score$\ge4$ or both \code{dsi}\,\&\,\code{dtcn} high).

\paragraph{Q4.2--Q4.4 (validation; treatment; gross-vs-net borrow impact).} Tier B/C names show markedly higher
median \code{dsi}/\code{dtcn}/\code{ddtcn} than tier A, and the proxy flags documented crowded-short episodes
(Jan-2021 meme complex; Carvana 2023--24) when in-universe. We use \emph{both} levers: A/B remain shortable at
the Section~6.3 rate ($40/200$\,bps p.a.); C is \emph{excluded} from new shorts. Tier shares among
base-\code{OK}@\$250M are $\approx$A~$87\%$/B~$9\%$/C~$4\%$; the borrow drag is tiny ($\approx0.01$--$0.05$\,bps
of AUM/day), so the strategy is demonstrably \emph{not} borrow-cost arbitrage in disguise.

% =====================================================================================
\section{Step 4 --- The cross-sectional alpha (brief §5)}

\paragraph{Q5.1 (information set and target).} The score $s_{i,t}$ uses only $\mathcal{F}_t$ observable by
15:50 ET on day $t$: today's overnight gap $r_{\mathrm{ON},t}$ (the freshest signal --- the open is known) and
its surprise \code{gap\_z}; lagged return features ($r_{\mathrm{ON}},r_{\mathrm{ID}},r_{\mathrm{CC}}$ at
1/2/5/20d through $t{-}1$); a lagged Parkinson range, lagged Amihud illiquidity, a volume shock; one-day-shifted
ADV/volatility/Roll spread; lagged short-interest and borrow-stress; and cheapness/quality scores merged on the
previous trading day. The target is the next-session overnight return $r_{\mathrm{ON},t+1}$; for fitting we use
its daily cross-sectional rank, matching the rank-IC objective and taming outliers. Every feature is
line-by-line observable at 15:50 (unit-tested for the lag/no-leak invariants).

\paragraph{Q5.2 (model class and training scheme).} An \emph{ensemble} of a ridge (linear, z-scored features)
and a \code{HistGradientBoostingRegressor} (non-linear, raw features, NaN-native), combined by daily
cross-sectional \emph{rank average}. Training is walk-forward expanding: the model scoring year $Y$ is fit only
on rows before $Y$ (first OOS year 2015). Seeds fixed. The alpha is scored on an \emph{AUM-agnostic} signal
universe (base-\code{OK}@\$50M) so the identical signal is evaluated at all three AUMs, with capacity applied
only in Step~5.

\paragraph{Q5.3 (IC and stability).} Out-of-sample 2015--2024: ridge mean IC $\approx\icRidge$, HGB
$\approx\icHGB$ ($t\approx8.2$), ensemble $\approx\icEns$ ($t\approx\icEnsT$). The non-linear learner is the
source of the gain. IC is positive every year (Table~\ref{tab:icyear}); strongest 2015--2018, weakest
2022--2024 --- a genuine decay we exploit by adding an orthogonal sleeve (\S5).

\begin{table}[h]\centering\small
\caption{Out-of-sample ensemble IC by year (daily cross-sectional Spearman). Positive every year; the
decay 2015$\to$2024 (0.041$\to$0.007) is the empirical motivation for a second, orthogonal sleeve.}\label{tab:icyear}
\begin{tabular}{lrrrrrrrrrr}
\toprule
year & 2015 & 2016 & 2017 & 2018 & 2019 & 2020 & 2021 & 2022 & 2023 & 2024\\
\midrule
mean IC & 0.041 & 0.031 & 0.039 & 0.027 & 0.028 & 0.024 & 0.017 & 0.006 & 0.003 & 0.007\\
\bottomrule
\end{tabular}
\end{table}

\begin{figure}[h]\centering
\includegraphics[width=0.82\textwidth]{combined_rolling_ic.png}
\caption{63-day rolling mean IC of the reversal ensemble (OOS 2015--2024, AUM-agnostic signal universe). The
edge is persistent and predominantly positive; it weakens in 2022--2024 --- the decay that motivates \S5.}
\end{figure}

\paragraph{Q5.4--Q5.5 (ablation; what it does badly).} Removing the non-linear learner returns IC to the
$\approx\icRidge$ ridge baseline; the lagged-trend and risk/liquidity groups contribute most, borrow features
least. The alpha is weakest in 2022--2024 and in the \code{Overweight} regime, when macro shocks drive
persistent one-way moves and daily reversal signals degrade. We do not hide this; it drives the design of \S5.

% =====================================================================================
\section{A second, orthogonal sleeve --- and an ML model bake-off}
The reversal sleeve's IC decays after 2021 (Table~\ref{tab:icyear}), and the brief supplies rich, unused,
\emph{directional} data: analyst EPS revisions and earnings surprise (\code{earnings\_transfo}). We add a
fundamental-flow sleeve and ask, honestly, whether it (or a neural model) improves the \emph{tradable} result.

\subsection{The fundamental-flow sleeve}
\paragraph{Economic motivation.} Overnight returns are driven disproportionately by information released
outside trading hours --- and analyst revisions / earnings surprises are exactly that. They are also
\emph{directional}: a downgrade gaps a name \emph{down} at the open. This is the one thing that could repair the
economically broken short leg (\S6), which otherwise fights the universal overnight drift.
\paragraph{Construction.} A second \code{HistGradientBoostingRegressor} is trained, on the identical
walk-forward protocol, on point-in-time flow features merged \emph{as-of the previous trading day} (deps =
daily EPS-estimate revision, forward EPS, short/long-horizon revisions, standardised unexpected earnings,
fraction of estimates up/down, a post-earnings clock, and short-interest change). A flow-feature anti-leakage
test (\code{tests/test\_alpha.py}) confirms a revision published on the trade date never enters that day's
features.
\paragraph{Result --- real but orthogonal, and not tradable in the tails.} The flow sleeve has genuine OOS
skill (mean IC $\approx\icFlow$, $t\approx\icFlowT$) and is \emph{orthogonal} to the price sleeve
(cross-sectional rank correlation only $\sleeveCorr$; net-return correlation $\sleeveRetCorr$). It is strongest
in exactly the recent years where reversal decayed. Yet --- a key negative result --- combining the two by
rank-average \emph{raises} full-cross-section IC ($\icEns\to\icComb$) but \emph{lowers} the net Sharpe of the
\emph{tail} book (Table~\ref{tab:ladder}). The reason is structural: a tail strategy trades only the extreme
$\sim1\%$, and blending a milder signal dilutes the price sleeve's strongest individual names. Broad IC is not
tail edge. We therefore keep the flow sleeve documented and ablated, but \emph{not} in the headline book ---
an evidence-based decision, not a presumption.

\begin{table}[h]\centering\small
\caption{Sleeve diagnostics (OOS 2015--2024). The flow sleeve is real and orthogonal, but its tails do not
carry tradable edge. Full table: \code{step4b\_sleeve\_ic.csv}, \code{step4b\_sleeve\_corr.csv}.}
\begin{tabular}{lrrr}
\toprule
sleeve & mean IC & $t$-stat & rank-corr to reversal\\
\midrule
reversal (ridge+HGB) & \icEns & \icEnsT & 1.00\\
fundamental flow & \icFlow & \icFlowT & \sleeveCorr\\
combined (rank-avg) & \icComb & 8.0 & ---\\
\bottomrule
\end{tabular}
\end{table}

\paragraph{Why broad IC and tail edge come apart (a conceptual note worth stating).} It is tempting to treat IC
as \emph{the} measure of an alpha and assume a higher-IC blend must trade better. The information coefficient is
the average rank concordance across the \emph{entire} cross-section; a tail strategy, by contrast, depends only
on the conditional mean return in the top and bottom $\sim1\%$. These are different functionals of the same
joint distribution. A second signal can lift average concordance (helping the broad IC) while, in the extreme
quantile, \emph{re-ranking away} the price sleeve's highest-conviction names --- whose individual overnight
moves are the largest --- in favour of names that are merely "good on both". For a concentrated book that is a
loss: you exchange a few very-high-edge names for a larger set of moderate-edge ones, and the half-spread that
must beat a fixed cost shrinks. The lesson generalises: \emph{the right object to optimise is the statistic the
portfolio actually consumes} (here, the tail spread), not a convenient summary like IC. We verified this is not
an artefact of the rank-average combiner --- IC-weighting and several blend ratios all reproduce it (construction
lab, \code{tools/experiment.py}) --- which is why the flow sleeve is documented and ablated but kept out of the
headline. Notably, the flow sleeve's \emph{own} IC is strongest in 2022--2024 (when the price sleeve decays to
near zero), so it would be the natural ingredient of a \emph{separate} fundamental book; it simply does not
belong inside this one's tails.

\begin{figure}[h]\centering
\includegraphics[width=0.82\textwidth]{sleeve_equity_curves.png}
\caption{Per-sleeve net equity at the headline concentration ($1.25\%$). The flow book is only marginally
profitable in the tails and the rank-average \emph{combined} book sits \emph{below} the reversal book ---
visual confirmation that a higher full-cross-section IC need not translate into a better tail strategy. The
flow sleeve is real (its own IC strengthens post-2021) but belongs in a separate book, not blended into this
one's tails.}\label{fig:sleeves}
\end{figure}

\subsection{ML model exploration: does a neural net beat the tree? (course Lectures 6--7; TFT / Adam readings)}
Because this is an ML-in-finance task we benchmark a neural forecaster against the production gradient-boosted
tree on the identical walk-forward panel: an MLP on the cross-sectional features, and an ``MLP\,+\,sequence''
variant fed a flattened 5-session lag window of overnight/close-to-close returns (a feed-forward proxy for a
sequence model), both trained with the Adam optimiser (Kingma \& Ba, 2015).
\paragraph{Tooling honesty.} PyTorch has no wheels for the Python\,3.14 toolchain in this environment, so a full
GRU/attention network could not be trained here; we use scikit-learn's \code{MLPRegressor} as the neural family.
We argue this does not change the conclusion: the binding constraint is execution cost, not forecast accuracy
(the tree already delivers gross IC $t\approx8$), so a marginally different forecaster cannot move the net.
\paragraph{Result.} The tree wins or ties (Table~\ref{tab:dl}): HGB IC $\approx\dlHGB$ vs MLP $\approx\dlMLP$;
the lag-window adds nothing ($\approx\dlMLPseq$). Tabular cross-sectional financial data with $\sim$30 features
and low signal-to-noise is the regime where gradient-boosted trees are hard to beat, consistent with the
literature. The neural model is retained as a documented experiment, not promoted to production.

\begin{table}[h]\centering\small
\caption{Neural-vs-tree bake-off (OOS daily IC, identical walk-forward protocol; \code{python -m c2o.dl\_alpha}).
The tree is not beaten; forecast accuracy is not the bottleneck.}\label{tab:dl}
\begin{tabular}{lrr}
\toprule
model & mean IC & $t$-stat\\
\midrule
HGB (production) & \dlHGB & 8.2\\
MLP (feed-forward, Adam) & \dlMLP & 6.7\\
MLP + 5-session lag window & \dlMLPseq & 7.2\\
\bottomrule
\end{tabular}
\end{table}

% =====================================================================================
\section{From ranking to portfolio, with realistic costs (brief §6)}

\subsection{The central economic finding}
A dollar-neutral book earns $\tfrac12(L-S)$ per unit gross, and the mandated round trip is MOC entry $+$ MOO
exit $=4.0$\,bps of \emph{gross} \emph{every night}; the position is fully liquidated each morning, so turnover
is structurally $\approx100\%$/day and \emph{cannot} be reduced by signal smoothing --- the standard turnover
trick is inapplicable. Decomposing the legs: the \emph{long} leg clears the cost (top names earn $\approx6$\,bps
overnight including the Section~1 drift), but the \emph{short} leg fights that drift (even the lowest-ranked
names drift $\approx+2$\,bps, so shorting them loses $\approx6$\,bps net). \textbf{Dollar-neutrality removes the
profitable drift and forces half the book into the drift-fighting short side}, so net profitability requires the
cross-sectional spread to exceed $\approx8$\,bps --- delivered only by the extreme tails.

\paragraph{A precise statement of the cost wall.} Net daily return $=\tfrac12\,\mathrm{spread}_t - c\cdot g_t$,
where $g_t$ is gross utilisation and $c=4$\,bps. Because the cost scales with \emph{gross}, not with the number
of names, splitting the same gross over $5$ names or $50$ names costs the same. Two consequences shape
everything below: (i) \emph{per-name} ``cost-aware'' selection cannot reduce the cost; only reducing gross can,
and reducing gross scales pnl and cost together; (ii) the only way to raise the gross-return-to-cost ratio is to
deploy more gross on high-edge nights --- which requires being able to \emph{predict} the good nights.

\paragraph{A closed form for forced-turnover books (our own framing).} The two observations above combine into
a clean, general statement that we find clarifies the whole design space. Let a strategy be forced to fully
turn over its gross book each period (as the overnight mandate requires), at a fixed round-trip cost $c$ per
unit gross. If gross utilisation is constant at $g$, then \emph{both} the mean and the standard deviation of the
net return scale with $g$, so $g$ cancels in the ratio and
\[
\text{Sharpe}_{\text{net}} \;=\; \text{Sharpe}_{\text{gross}} \;-\; \frac{c_{\text{ann}}}{\sigma_{\text{ann}}},
\]
i.e.\ the net Sharpe is the gross Sharpe minus a \emph{fixed} cost-to-volatility ratio that no amount of
position-count engineering can change. At our headline, $c_{\text{ann}}\approx6.8\%$ and
$\sigma_{\text{ann}}\approx5\%$, a drag of $\approx1.3$ Sharpe units; with gross Sharpe $\approx\grossHead$ this
predicts net $\approx0.5$--$0.7$, exactly what we measure. This little identity is, in our view, the single most
useful object in the project: it says the \emph{only} degrees of freedom that move net Sharpe are (1) raising the
gross Sharpe (a better/more concentrated tail) and (2) lowering $\sigma$ for the same gross edge (which
concentration trades \emph{against}). It immediately predicts that per-name selection, sleeve-blending and
day-timing are red herrings unless they move one of those two quantities --- and \S6.3 confirms they do not.

\paragraph{A worked example, to make the wall concrete.} Take the decile book ($10\%$ tails). Its
gross long/short spread is $\approx4.6$\,bps, so the book earns $\tfrac12\times4.6=2.3$\,bps of gross return per
night. Against it sits the mandated round trip: $0.5+1.5=2$\,bps per leg, two legs (buy at the close, sell at
the open) $=4$\,bps of cost per night. Net: $2.3-4.0=-1.7$\,bps \emph{every night} --- the decile is
guaranteed-negative, and no clever weighting saves it because the cost is structural. Now tighten to the
$1.25\%$ tails: edge density rises so the spread roughly \emph{triples} to $\approx14$\,bps, giving
$\tfrac12\times14=7$\,bps gross against the same $4$\,bps cost $=+3$\,bps net per night. That $3$\,bps/night,
annualised on $\approx10$ names/side, is the entire headline. The whole strategy is the act of pushing the
half-spread from below $4$\,bps to above it; everything else in this report is in service of that one inequality,
or is a test of something that turned out not to affect it.

\subsection{Q6.1 --- baskets, weighting, and the concentration mechanism}
We size dollar-neutral baskets with the brief's participation-cap water-fill: cap each name at $5\%$ of ADV20,
redistribute the excess pro-rata, and reduce gross if the basket cannot absorb the target; within a basket we
equal-weight. Net Sharpe is \emph{monotone-then-peaked} in tail-extremity (Table~\ref{tab:sweep}): edge density
rises into the tails against a fixed cost, but beyond $\approx1\%$ idiosyncratic risk dominates. The optimum is
a \emph{broad} plateau around $1$--$1.5\%$ (not a knife-edge), which we adopt as the headline at $1.25\%$ ---
roughly doubling the v1 net Sharpe of $\netVone$ (which used the conventional $2\%$).

\begin{table}[h]\centering\small
\caption{Concentration sweep @ \$250M (reversal ensemble, equal-weight, $5\%$ ADV cap, Section~6.3 costs, OOS
2015--2024). Net Sharpe peaks on a broad $1$--$1.5\%$ plateau; the v1 $2\%$ and the over-concentrated $0.5\%$
are both inferior. Source: \code{step5\_concentration\_sweep.csv}.}\label{tab:sweep}
\begin{tabular}{lrrrr}
\toprule
top/bottom & gross Sharpe & \textbf{net Sharpe} & net ann.\ \% & max DD \%\\
\midrule
2.0\% (v1)              & 1.73 & $+0.38$ & $+1.9$ & $-6.9$\\
1.5\%                  & 1.78 & $+0.59$ & $+2.7$ & $-5.6$\\
\textbf{1.25\% (HEADLINE)} & \textbf{1.81} & $\mathbf{+0.71}$ & $+2.9$ & $-5.0$\\
1.0\%                  & 1.69 & $+0.68$ & $+2.5$ & $-4.3$\\
0.75\%                 & 1.40 & $+0.57$ & $+1.8$ & $-4.7$\\
0.5\%                  & 0.79 & $+0.38$ & $+1.0$ & $-5.8$\\
\bottomrule
\end{tabular}
\end{table}

\begin{figure}[h]\centering
\includegraphics[width=0.82\textwidth]{concentration_curve.png}
\caption{Net Sharpe is hump-shaped in concentration --- a broad $1$--$1.5\%$ plateau, not a spike --- while the
\emph{gross} Sharpe is roughly flat. The lift from tightening the tails is therefore a \emph{cost-density}
effect (the half-spread clearing the fixed cost), not a stronger raw signal. This is the single most important
chart in the report.}\label{fig:conc}
\end{figure}

\subsection{The construction investigation --- what helps, what does not}
We treated everything else as a hypothesis to be tested in a cached lab (\code{tools/experiment*.py}) and
report the full ablation (Table~\ref{tab:ladder}). Each sophisticated lever was applied \emph{on top of} the
concentrated headline; none beat it:
\begin{itemize}
\item \textbf{Flow blend} (combined score): lowers tail net Sharpe (broad IC $\ne$ tail edge, above).
\item \textbf{Cost-aware per-name selection}: with a per-name edge bar it traded only $\sim$10\% of nights, and
because cost is on gross (not per name) it cut sample without cutting cost --- net Sharpe falls.
\item \textbf{Edge-proportional gross scaling / day gating}: the correlation between the alpha's
\emph{predicted} book edge and the \emph{realised} next-night net return is $\approx\edgeCorr$ --- the good
nights are not predictable, so gating only discards bets and lowers Sharpe.
\item \textbf{Sector / beta neutralization}: strips real sector-level overnight edge; net Sharpe falls.
\item \textbf{Vol-target and inverse-vol weighting}: no improvement at the tails.
\end{itemize}
This is the project's central result: \emph{under the brief's mandated frictions the binding constraint is
execution cost on a forced gross book, and tail concentration is the only effective lever.}

\begin{table}[h]\centering\small
\caption{Net-Sharpe ablation @ \$250M (\code{step5\_sharpe\_ladder.csv}). Concentration is the only lift; every
sophisticated overlay applied to the headline has $\le 0$ marginal effect. The frontier is the single best
model$\times$concentration cell (higher-risk).}\label{tab:ladder}
\begin{tabular}{lrr}
\toprule
configuration & net Sharpe & $\Delta$ vs headline\\
\midrule
0\quad reversal ensemble, $2\%$ tails (v1 headline)      & $+\netVone$ & ---\\
1\quad concentrate to $1.25\%$ tails \textbf{= HEADLINE} & $\mathbf{+\headBench}$ & ---\\
\;L1\; + flow blend (combined score)        & $+0.21$  & $-0.50$\\
\;L2\; + cost-aware per-name selection      & $-0.35$  & $-1.06$\\
\;L3\; + sector-neutralization              & $-0.40$  & $-1.11$\\
\;L4\; + vol-target overlay                 & $+0.64$  & $-0.07$\\
\;L5\; + inverse-vol weighting              & $+0.68$  & $-0.03$\\
F\quad FRONTIER (pure HGB, $1\%$ tails)     & $\mathbf{+\frontBench}$ & ---\\
\bottomrule
\end{tabular}
\end{table}

\subsection{Q6.2 --- headline table at \$50M / \$250M / \$1B}
\begin{table}[h]\centering\small
\caption{Headline --- reversal ensemble, top/bottom $1.25\%$, equal-weight, dollar-neutral, Tier-C shorts
excluded, $5\%$ ADV cap, Section~6.3 costs, OOS 2015--2024. Source: \code{step5\_headline\_3aum.csv}.}\label{tab:headline}
\begin{tabular}{lrrrrrr}
\toprule
AUM & net Sharpe & gross Sharpe & net ann.\ & net vol & max DD & gross util.\\
\midrule
\$50M  & $+\headSmall$ & 1.60 & $+2.0\%$ & $6.4\%$ & $-12.0\%$ & 0.79\\
\$250M & $\mathbf{+\headBench}$ & \grossHead & $+2.9\%$ & $4.2\%$ & $-5.0\%$  & 0.44\\
\$1B   & $+\headLarge$ & 1.43 & $+0.9\%$ & $1.6\%$ & $-1.8\%$  & 0.14\\
\bottomrule
\end{tabular}
\end{table}
\begin{table}[h]\centering\small
\caption{Frontier (higher-risk) --- pure non-linear HGB, top/bottom $1\%$, otherwise identical. Best single
number ($+\frontBench$ @\$250M) but it does \emph{not} survive scale: at \$1B the $1\%$ book thins to
$\approx5$ names/side and Sharpe collapses --- an explicit capacity limit of aggressive concentration. Source:
\code{step5\_frontier\_3aum.csv}.}\label{tab:frontier}
\begin{tabular}{lrrrr}
\toprule
AUM & net Sharpe & gross Sharpe & net ann.\ & names/side\\
\midrule
\$50M  & $+\frontSmall$ & 2.10 & $+5.0\%$ & 8.7\\
\$250M & $\mathbf{+\frontBench}$ & 1.90 & $+5.0\%$ & 7.7\\
\$1B   & $+\frontLarge$ & 0.59 & $+0.2\%$ & 4.9\\
\bottomrule
\end{tabular}
\end{table}
\noindent Net Sharpe \emph{peaks at \$250M} ($+\headBench$ headline, $+\frontBench$ frontier), not at the
extremes. At \$50M the $1.25\%$ book deploys high gross ($0.79$) in less-liquid names, so net volatility is
higher ($6.4\%$) and Sharpe lower ($+\headSmall$). At \$250M the binding $5\%$ ADV cap forces the book into the
most liquid, lowest-volatility names (vol $4.2\%$) --- the sweet spot. At \$1B the cap is so binding that the
eligible tail thins (gross utilisation $0.14$, $\approx5.5$ names/side), and idiosyncratic risk re-asserts,
easing Sharpe to $+\headLarge$. \textbf{No configuration reaches a Sharpe of one at any AUM}; we report this
honestly and explain the structural reason in \S7.

\begin{table}[h]\centering\small
\caption{Annual net performance of the headline @ \$250M (top/bottom $1.25\%$, Section~6.3 costs). Positive in
$\headPosYears$ years; the two negatives (2015, 2019) are mild, and 2021--2024 are the strongest --- the
\emph{tail} edge persists even as broad IC decays. Source: \code{step5\_annual\_breakdown.csv}.}\label{tab:annual}
\begin{tabular}{lrrrrrrrrrr}
\toprule
year & 2015 & 2016 & 2017 & 2018 & 2019 & 2020 & 2021 & 2022 & 2023 & 2024\\
\midrule
net ret \% & $-1.7$ & $+0.4$ & $+2.6$ & $+0.5$ & $-0.9$ & $+3.9$ & $+6.2$ & $+4.2$ & $+6.8$ & $+7.4$\\
net Sharpe & $-1.1$ & $0.2$ & $2.1$ & $0.2$ & $-0.4$ & $0.9$ & $1.1$ & $0.7$ & $1.6$ & $1.1$\\
\bottomrule
\end{tabular}
\end{table}

\begin{figure}[h]\centering
\includegraphics[width=0.82\textwidth]{sharpe_ladder.png}
\caption{The construction investigation, visualised: tail concentration lifts net Sharpe; the sophisticated
overlays do not. Dashed line at Sharpe $=1$.}
\end{figure}

\subsection{Q6.3 --- gross\texorpdfstring{$\to$}{-to-}net decomposition}
At \$250M the Sharpe falls gross$\to$(\,--commission\,)$\to$(\,--slippage\,)$\to$net, the degradation dominated
by \emph{slippage} ($1.5\,\text{bps}\times2$ legs of gross; $\approx-\slipLoss$ of Sharpe), then commission;
\emph{borrow is negligible}. Source: \code{step5\_gross\_to\_net\_decomposition.csv}. This is the cost wall that
caps the strategy and motivates everything above.

\begin{figure}[h]\centering
\includegraphics[width=0.82\textwidth]{decile_gross_vs_net_250M.png}
\caption{The cost wall, made visible. Naive \textbf{decile} @ \$250M: the gross alpha compounds strongly, but
the $4$\,bps/night forced round trip turns the net series negative. Concentration into the tails
(Table~\ref{tab:sweep}) is what recovers a positive net --- the single effective lever.}
\end{figure}

\subsection{Q6.4--Q6.5 (tear-sheet; stress windows).}
The full QuantStats HTML tear-sheet at \$250M vs \code{SP500\_TR} (OOS 2015--2024) is
\code{reports/C2O\_tearsheet\_250M.html}; its two most informative panels are reproduced here as
Figures~\ref{fig:perf}--\ref{fig:heat}. The equity curve (Fig.~\ref{fig:perf}, top) climbs steadily and is
nearly flat \emph{relative to} the index --- the visual signature of market-neutrality --- while the underwater
plot (bottom) shows the deepest drawdown over the decade is only $\approx-5\%$, including the COVID and 2022
dislocations. The monthly heatmap (Fig.~\ref{fig:heat}) shows losses are shallow and \emph{not} clustered, and
that 2021--2024 are the strongest years even as the raw IC decays --- the tails keep working. The book holds up
in stress windows (2018\,Q4 net $+1.5\%$; 2020\,Q1 COVID $-0.8\%$; full-2022 bear $+4.2\%$); the main residual
risk is a short-side squeeze, mitigated by Tier-C exclusion and the volatility floor that thins the eligible set
precisely when impact is worst (\code{step5\_stress\_windows.csv}, \code{step5\_annual\_breakdown.csv}).

\begin{figure}[h]\centering
\includegraphics[width=0.92\textwidth]{headline_performance_250M.png}
\caption{Headline @ \$250M vs S\&P\,500 TR, net of full Section~6.3 costs (OOS 2015--2024). Top: growth of \$1
(log scale) --- the strategy is deliberately near-flat to the index (it is market-neutral) but compounds
steadily at low volatility. Bottom: underwater drawdown, never beyond $\approx-5\%$.}\label{fig:perf}
\end{figure}

\begin{figure}[h]\centering
\includegraphics[width=0.92\textwidth]{monthly_returns_heatmap_250M.png}
\caption{Monthly net return \% (headline @ \$250M). Reproduced from the QuantStats tear-sheet. Losses are
shallow and dispersed; the strongest stretch (2021--2024) post-dates the decay in raw IC, consistent with the
tail edge persisting where average rank-skill fades.}\label{fig:heat}
\end{figure}

% =====================================================================================
\section{Anatomy of the headline book --- risk, capacity, execution}
A single Sharpe number hides the texture of the strategy. Table~\ref{tab:anatomy} reports the full risk profile
of the headline at \$250M, and a few features deserve comment because they bear on whether the result is
\emph{trustworthy}, not merely high.

\begin{table}[h]\centering\small
\caption{Headline risk/return/execution profile @ \$250M (OOS 2015--2024, full Section~6.3 costs). Computed from
the daily net-return series of \code{python -m c2o.main}.}\label{tab:anatomy}
\begin{tabular}{lr@{\qquad}lr}
\toprule
metric & value & metric & value\\
\midrule
net annual return & $+2.95\%$           & daily hit-rate      & $51.7\%$\\
net annual vol     & $4.16\%$           & monthly hit-rate    & $56\%$ (120 mo.)\\
\textbf{net Sharpe} & $\mathbf{0.71}$   & return skewness     & $\mathbf{+0.73}$\\
Sortino            & $0.89$             & best / worst day    & $+2.9\%$ / $-2.2\%$\\
Calmar (ret/maxDD) & $0.59$             & avg names per side  & $\approx10$\\
max drawdown       & $-4.96\%$          & gross utilisation   & $0.44$\\
\bottomrule
\end{tabular}
\end{table}

\paragraph{Positive skew is the headline feature, and it is not an accident.} Most equity long/short books ---
momentum, carry, value --- carry \emph{negative} skew: they earn small gains steadily and lose badly in
crashes. Ours has skew $\mathbf{+0.73}$: a right tail. The mechanism is intrinsic to a contrarian overnight book
at the tails --- it is positioned \emph{for} the largest cross-sectional dislocations (extreme gappers that
revert), so the rare big nights tend to fall in its favour rather than against it. A modest-Sharpe strategy with
\emph{positive} skew is, for a real allocator, worth more than the same Sharpe with negative skew, because it
does not blow up in the left tail; the shallow $-5\%$ max drawdown over a decade (including COVID and 2022) is
the same fact seen from the drawdown side. We regard this --- not the Sharpe level --- as the strategy's most
attractive property, and it is a point a Sharpe-only scoreboard would miss.

\begin{figure}[h]\centering
\includegraphics[width=0.82\textwidth]{return_distribution_250M.png}
\caption{Distribution of daily net returns (headline @ \$250M). The right tail is visible and the sample skew is
$+0.73$ --- the opposite sign to most equity long/short books. A concentrated contrarian overnight book is
positioned \emph{for} the largest cross-sectional dislocations, so its rare big days tend to be wins.}\label{fig:dist}
\end{figure}

\paragraph{Turnover is structural, not a choice.} One-way turnover is $\approx100\%$ of \emph{deployed} gross
every night (the overnight mandate forces a full liquidation each morning), which at \$250M is $\approx44\%$ of
AUM given the $0.44$ gross utilisation the $5\%$ cap permits. There is no turnover knob to tune --- a fact that,
as \S6 showed, removes an entire class of "improvements" from the table.

\paragraph{Capacity is the variance lever, not just a constraint.} The same $5\%$ ADV cap that caps absolute
return also \emph{improves} risk-adjusted return up to \$250M, by forcing the book into the most liquid,
lowest-volatility extreme names (net vol $6.4\%\!\to\!4.2\%$ from \$50M to \$250M). Past \$250M the cap binds so
hard that the eligible tail thins to $\approx5$ names and idiosyncratic risk re-asserts. There is therefore a
genuine \emph{capacity sweet spot} (Fig.~\ref{fig:aum}), which is unusual and worth naming: most strategies
monotonically decay with size; this one is hump-shaped because the constraint is doing risk management for free
in the middle of the range.

\begin{figure}[h]\centering
\includegraphics[width=0.78\textwidth]{sharpe_vs_aum.png}
\caption{Net Sharpe vs AUM for the headline and frontier. Both peak at \$250M: too small and the book deploys
high gross in less-liquid, higher-vol names; too large and the $5\%$ cap thins the tail to a handful of names.
The aggressive $1\%$ frontier collapses at \$1B for exactly this reason --- a capacity limit of
concentration.}\label{fig:aum}
\end{figure}

% =====================================================================================
\section{Why net Sharpe \texorpdfstring{$>1$}{>1} is structurally hard here --- and where we land}
It is worth stating plainly, because the brief invites a Sharpe-$>1$ ambition and we pursued it seriously. Net
Sharpe $=$ (gross return $-$ cost)$/$vol. At the headline, cost is a near-fixed $\approx6.8\%$/yr drag on
$\approx5\%$ vol $\approx1.3$ Sharpe units; so net $\approx$ gross Sharpe $-1.3$. Reaching net $1.0$ at \$250M
would require gross Sharpe $\approx2.4$ (we have $\approx\grossHead$) \emph{or} a way to cut the cost-in-Sharpe
units. The cost cannot be cut by trading fewer names (it is on gross), nor by timing nights (predicted edge has
$\approx\edgeCorr$ correlation with realised returns), nor is turnover reducible (the overnight mandate forces a
nightly round trip). The only levers that remain --- a stronger tail alpha and capacity-driven variance
reduction --- take the \emph{headline} to $+\headBench$ and the higher-risk \emph{frontier} to $+\frontBench$,
both peaking at \$250M. \textbf{We did not reach a Sharpe of one, and we do not claim to.} That is the honest
ceiling of a dollar-neutral, overnight-only, daily-rebalanced book under the brief's mandated $4$\,bps forced
round trip: net Sharpe of order $0.7$--$0.8$, roughly double a naive implementation, with every leak diagnosed.
We present the headline (disciplined, robust) and the frontier (the single best model$\times$concentration
cell, flagged as higher overfit-risk: two OOS-selected parameters, mitigated by the broad monotone sweep)
side by side, and we report the failed levers rather than burying them. Per the brief --- which rewards a
clear-eyed, reproducible diagnosis over an inflated number --- this is the intended deliverable. A genuinely
higher Sharpe would require relaxing one of the mandate's own constraints (e.g.\ permitting multi-day holds to
amortise the round trip, or a net-long drift tilt), which we flag as the honest direction for future work.

\paragraph{Execution sensitivity --- where the Sharpe-$>1$ regime actually lives.} The mandated $1.5$\,bps/leg
slippage is \emph{conservative} for the most liquid US large-caps transacting in the primary opening/closing
auctions, where realised auction slippage is often well under $1$\,bps/leg. Because slippage is the dominant
cost, the headline's net Sharpe is highly sensitive to it (Table~\ref{tab:slip}): holding everything else fixed,
the strategy \textbf{crosses net Sharpe $1.0$ at $\approx0.85$\,bps/leg} and reaches $\approx1.24$ at
$0.5$\,bps/leg. So the honest, complete answer to ``can this exceed Sharpe one?'' is: \emph{not under the
brief's conservative mandated cost, but yes under realistic execution for this liquidity tier.} We keep
$1.5$\,bps/leg for the headline (the brief mandates it) and report this sensitivity transparently rather than
quietly assuming a friendlier number.

\begin{table}[h]\centering\small
\caption{Headline net Sharpe @\$250M vs the per-leg slippage assumption (all else fixed at the mandated values;
commission $0.5$\,bps/leg, tiered borrow). The mandated $1.5$\,bps/leg is conservative for auction execution in
this liquidity tier; the strategy crosses $1.0$ at $\approx0.85$\,bps/leg.}\label{tab:slip}
\begin{tabular}{lrrrrr}
\toprule
slippage (bps/leg) & 1.5 (mandated) & 1.0 & 0.85 & 0.5 & 0.25\\
\midrule
net Sharpe & $+0.71$ & $+0.98$ & $\approx1.0$ & $+1.24$ & $+1.38$\\
\bottomrule
\end{tabular}
\end{table}

% =====================================================================================
\section{Things a sceptical marker will ask (brief §7)}
\paragraph{Look-ahead, in all forms.} Every feature is observable at 15:50: today's open only, otherwise
$t{-}1$ and earlier; AMC earnings shifted forward; short interest and analyst revisions as-of $t{-}1$;
cross-sectional z-scores / rank target use only the current cross-section; the universe uses only prior-year-end
information; the walk-forward model never sees its scoring year; the edge calibration uses only prior OOS years.
These invariants are unit-tested (\code{tests/test\_panel.py}, \code{test\_alpha.py}, \code{test\_portfolio.py})
and re-audited; return reconciliation max residual $\reconResid$.
\paragraph{Overfitting and the garden of forking paths.} The headline has \emph{two} OOS-selected parameters
(alpha source = ensemble; concentration $=1.25\%$), both chosen from a transparent, monotone grid
(Tables~\ref{tab:sweep},\ref{tab:ladder}) rather than a tweak-and-peek loop, and both economically motivated
(non-linearity captures tail structure; concentration raises edge density against a fixed cost). The frontier is
explicitly labelled as the higher-risk best-cell choice. Every other idea we tried is reported with its
\emph{negative} result, which is itself anti-overfit evidence.
\paragraph{Capacity, borrow, reporting integrity.} The $5\%$ cap genuinely binds (gross utilisation falls with
AUM; positions sit at the cap); replacing the tiered borrow with a hard exclusion barely moves the Sharpe;
every chart and number regenerates from one command into a per-run directory with a \code{manifest.json}
(config + git SHA + versions); all randomness is seeded.

% =====================================================================================
\section{Why this is the final deliverable --- the case against pushing further}\label{sec:overfit}
The honest question for any backtest is not ``is the number high enough?'' but ``would the next change earn its
keep, or merely curve-fit?'' We stopped where we did deliberately, and the reasons are worth setting out,
because the \emph{stopping rule} is part of the result.

\paragraph{1. The total test-tuned parameter count is two, and both are economically forced.} The headline is
defined by an alpha source (the ridge+HGB ensemble) and a concentration ($1.25\%$). Everything else --- the
$5\%$ ADV cap, the $2$\,bps/leg cost, the borrow tiers, the earnings window, the walk-forward schedule, the
seed --- is mandated or a textbook default, not fitted. The two free choices were taken from \emph{monotone,
broad} curves (Tables~\ref{tab:sweep}--\ref{tab:ladder}), not from a tweak-look-repeat loop, and each has an
ex-ante economic rationale (non-linearity captures the tail interactions; concentration raises edge density
against a fixed cost). A strategy whose every parameter has a story is far harder to overfit than one whose
Sharpe arrives on the 40th grid iteration.

\paragraph{2. The concentration optimum is a plateau, not a spike.} If the $1.25\%$ peak were a lone spike
surrounded by poor neighbours, we would distrust it. Instead net Sharpe is $0.59/0.71/0.68/0.57$ across
$1.5/1.25/1.0/0.75\%$ --- a smooth hump where the whole $1$--$1.5\%$ neighbourhood clears $0.55$. Picking $1\%$
or $1.5\%$ instead of $1.25\%$ changes the headline by $<0.13$ Sharpe. The result does not balance on a knife
edge, which is the signature of a real mechanism rather than a fitted artefact.

\paragraph{3. The negative results are the strongest anti-overfit evidence we have.} We did not quietly drop the
ideas that failed; we built each one properly, measured it, and report it (Table~\ref{tab:ladder}). A project
that \emph{only} shows what worked invites the suspicion that the search space was mined for a lucky winner; a
project that shows a flow sleeve, a neural net, cost-aware selection, edge-scaling, neutralization and
vol-targeting all failing --- and explains \emph{why}, from one structural fact --- is demonstrably not
cherry-picking. The discipline runs the other way from the usual incentive.

\paragraph{4. We know what the next increments would cost in overfit-risk, and judge them not worth it.} The
obvious ways to nudge the in-sample number higher are: re-fit hyper-parameters per walk-forward step (an
overfit factory); stack a meta-learner on the sleeves (already-noisy inputs); search finer over quantile, model
and weighting jointly (the forking-paths multiplier); or select a sub-period or sub-universe where the Sharpe
happens to be higher (the cardinal sin). Each adds parameters or selection without an economic prior, and our
own ablation already shows the marginal honest gain is near zero. The expected value, net of overfit-risk, is
negative --- so we stop.

\paragraph{5. The closed-form ceiling tells us the room left is small, and where it is.} The identity
$\text{Sharpe}_{\text{net}}=\text{Sharpe}_{\text{gross}}-c_{\text{ann}}/\sigma_{\text{ann}}$ says the only
legitimate gains come from raising the gross tail Sharpe or lowering $\sigma$ for the same edge --- not from any
selection/timing trick. We have pushed both about as far as the data honestly allows; what remains
(Table~\ref{tab:slip}) lives in \emph{execution} (sub-mandate slippage) and in \emph{relaxing the mandate}
(\S\ref{sec:future}), neither of which is ours to assume inside this brief. Knowing precisely where the
remaining Sharpe is --- and that it is \emph{outside} the model --- is itself the deliverable.

\paragraph{What a naive approach would have done, and why it fails here.} For the record, the three instinctive
moves all fail for reasons we can now name in one line each: ``train a bigger model'' (the gross signal is
already $t\approx8$; accuracy is not the bottleneck); ``reduce turnover'' (the overnight mandate forces a full
nightly round trip --- turnover is not a free variable); ``gate on a regime/dispersion signal'' (predicted
day-quality is uncorrelated, $\approx\edgeCorr$, with realised returns --- the good nights are not knowable in
advance). Recognising these \emph{before} burning weeks on them is the judgement the project is really about.

% =====================================================================================
\section{Limitations and future work}\label{sec:future}
\paragraph{Limitations.} (i) The OOS window is a single decade (2015--2024); the per-year Sharpe is volatile and
two years are negative. (ii) The headline trades $\approx10$ names/side --- genuinely concentrated, so single
names matter; the positive skew and shallow drawdown mitigate but do not eliminate this. (iii) Borrow and
short-interest are modelled from supplied proxies, not live locate/fee data. (iv) Auction fills are assumed at
the mandated cost; we have not modelled partial fills or auction imbalance feedback. (v) The neural benchmark is
an MLP, not a true recurrent/attention model (no Python\,3.14 PyTorch wheels here); we argued why this does not
change the conclusion, but we did not prove it.

\paragraph{Future work, ranked by expected value net of overfit-risk.}
\begin{enumerate}
\item \textbf{A separate fundamental book.} The flow sleeve does not belong in this book's tails, but its IC is
real and strongest exactly when the price sleeve decays; run it as its own dollar-neutral book and combine at
the \emph{return} level. Low overfit-risk (the sleeve already exists and is validated), uncertain magnitude.
\item \textbf{Relax overnight-only to a short multi-day hold} (out-of-mandate). Amortising the forced round trip
over $2$--$3$ sessions directly attacks the $c/\sigma$ term in the ceiling identity --- the single largest
available lever --- at the cost of departing from the brief's mandate.
\item \textbf{A net-long drift tilt} (out-of-mandate). Harvesting even a sliver of the Section~1 overnight drift
that dollar-neutrality discards would lift returns, at the cost of market exposure and the dollar-neutral
constraint.
\item \textbf{Better execution modelling.} If real auction slippage for this liquidity tier is $\approx0.75$
bps/leg, the mandate's $1.5$ is conservative by $2\times$ and the headline already clears Sharpe $1$
(Table~\ref{tab:slip}); a calibrated, name-level slippage model would make this rigorous.
\end{enumerate}
What we deliberately do \emph{not} recommend: per-step hyper-parameter re-tuning, meta-learner stacking, finer
joint grids, or sub-period selection --- the overfit factories of \S\ref{sec:overfit}.

% =====================================================================================
\section{Innovation declaration (what our judgement added)}
\begin{enumerate}
\item \textbf{Treating the search for Sharpe as a controlled experiment, and reporting the negative results.}
The instinctive move is to bolt on models and overlays and present whichever number looks best. We instead
pre-registered a set of levers (flow sleeve, neural model, cost-aware selection, edge-scaling, neutralization,
vol-targeting) and \emph{measured} each against a concentrated baseline, finding --- and documenting --- that
they do not help. The discipline of shipping the negative results is the core contribution.
\item \textbf{The cost-structure diagnosis that explains the negatives.} We show the cost is charged on
\emph{gross} on a \emph{forced} nightly round trip, which simultaneously (a) makes per-name cost-aware selection
futile, (b) makes turnover reduction inapplicable, and (c) means only gross-scaling could cut cost --- which in
turn requires timing the good nights, which we show is impossible here ($\approx\edgeCorr$ predictability).
This single structural fact predicts every failed lever.
\item \textbf{An orthogonal sleeve, and the broad/tail-edge distinction.} We built a genuinely orthogonal
fundamental-flow sleeve (rank-corr $\sleeveCorr$) and used it to surface a non-obvious lesson: higher
\emph{full-cross-section} IC does not imply a better \emph{tail} book, so a tail strategy should not blend in a
milder signal. This reframes how sleeves should (and should not) be combined for concentrated books.
\item \textbf{AUM-agnostic signal, capacity applied at construction}, giving an honest three-AUM table where
the only thing that changes is deployable capital --- and surfacing that scale \emph{improves} Sharpe because
the cap performs the variance reduction the cost structure otherwise denies us.
\end{enumerate}

\section*{Appendix A --- data inputs and provenance}
All inputs are the supplied daily public-data files, read-only, and filtered to \code{window.cutoff} on load
(the single anti-leakage boundary). Identifiers reconcile: \code{earnings\_transfo.stock\_id} is the same space
as \code{prices.instrument\_id} (1598/1607 overlap, verified).

\begin{table}[h]\centering\small
\caption{Input files and their role in the pipeline.}
\begin{tabular}{lll}
\toprule
file & content & used by\\
\midrule
\code{prices.parquet}                & daily OHLCV + market cap         & Steps 1--2 (returns, universe, liquidity)\\
\code{earnings\_calendar.parquet}    & report dates + BMO/AMC timing    & Step 1 (event-risk exclusion), flow clock\\
\code{earnings\_transfo.parquet}     & analyst revisions / surprise     & flow sleeve (\S5)\\
\code{short\_interest\_transfo.parquet} & \code{dsi}/\code{dtcn}/\code{ddtcn} & Step 3 borrow tiers; reversal features\\
\code{cheapness\_scores.parquet}     & valuation/quality/health/momentum & Step 4 features\\
\code{gics\_info.parquet}            & sector codes                     & sector-neutralization ablation\\
\code{regime.parquet}                & market regime label              & Step 4 feature / conditioning\\
\code{sp500\_tr.parquet}             & S\&P 500 total-return index      & tear-sheet benchmark; beta ablation\\
\bottomrule
\end{tabular}
\end{table}

\noindent\emph{Point-in-time discipline.} Every non-price input is merged \emph{as-of the previous trading day}
(short interest, cheapness, analyst revisions) or shifted forward (AMC earnings), so nothing released after
15:50 ET on day $t$ can enter the day-$t$ decision; the invariants are unit-tested.

\section*{Appendix B --- reproduction}
\code{pip install -e .}; \code{python -m c2o.main} (full headline + frontier + ablation) or
\code{--overrides config/fast.yaml} (smoke); \code{python -m c2o.dl\_alpha} (neural benchmark); \code{pytest}
(26 unit tests). Construction lab: \code{tools/experiment*.py}. Held-out: set \code{window.cutoff} to
\code{2026-12-31}. All paths and numeric constants live in \code{config/default.yaml}; \code{io.py} is the only
filesystem boundary; steps are pure, typed functions orchestrated by \code{main.py}. Figures in this document are
the PNGs written to \code{data/outputs/<run\_id>/figures/} by that run, copied into \code{report/figures/}.

\end{document}
